
Our four-stage validation pipeline tests proxy metrics through increasing rigor. Stage 1 (measurement reliability, validated in this work) filters noisy metrics before optimization. Stages 2-4 (conditional independence, cross-domain generalization, optimization constraints) are designed for future validation. This section presents Stage 1 in detail.

\subsubsection{Four-Stage Validation Framework}

\textbf{Stage 1: Measurement Reliability} tests whether a proxy exhibits stable, discriminative, generalizable measurements. Three criteria must all pass:

\begin{enumerate}
\item \textbf{Intra-implementation stability (CV $\leq 5$\%):} Measuring the same code solution five times yields low variance, indicating the metric is not dominated by measurement noise.
\item \textbf{Inter-complexity-class separation (Cohen's d $\geq 0.8)$:} The metric distinguishes algorithmic complexity classes (O(n) vs O($n^{2})$ solutions) with large effect size.
\item \textbf{Cross-platform generalization (Spearman $\rho \geq 0.8)$:} Rank ordering of solutions remains consistent across hardware platforms.
\end{enumerate}

Proxies failing any criterion proceed no further. The scoped gate design allows $\geq 1$ proxy to validate for progression.

The CV $\leq 5$\% threshold comes from psychometric reliability standards for continuous measures. Cohen's d $\geq 0.8$ ensures proxies distinguish meaningful quality differences (Cohen's ''large effect'' convention). Spearman $\rho \geq 0.8$ ensures cross-platform portability. These thresholds are provisional and domain-specific validation may adjust them.

\textbf{Stage 2: Conditional Independence} (designed for future validation in h-e2) tests whether proxies explain behavioral outcome variance after controlling for execution correctness through hierarchical regression. If $\Delta R^2 \geq 0.03$ and effects persist within the perfect-execution stratum, the proxy captures non-redundant information.

\textbf{Stage 3: Cross-Domain Generalization} (future work) validates that proxy effects generalize across repositories using leave-cluster-out cross-validation. Success requires prediction $R^2$ drops <50\% and ensemble disagreement $\leq 20$\%.

\textbf{Stage 4: Optimization Constraints} (future work) tests per-task execution safety during RL training through per-task monitoring enforcing $\leq 5$\% regression per problem.

\subsubsection{Stage 1 Implementation: Measurement Reliability}

\#\#\#\# Three-Proxy System

\textbf{Proxy 1: CodeBLEU (Structural Similarity).} CodeBLEU combines four sub-metrics: n-gram match, weighted n-gram match, AST match, and dataflow match. We use the \texttt{k4black/codebleu} implementation with equal weighting (0.25, 0.25, 0.25, 0.25). CodeBLEU is a deterministic computation---AST parsing and dataflow analysis have no randomness.

\textbf{Proxy 2: Runtime Efficiency (Algorithmic Performance).} Following the CPU time equation methodology (Patterson \& Hennessy), we measure efficiency via CPU instruction count rather than wall-clock execution time. The rationale: CPU time = [Instruction Count] $\times$ [CPI] $\times$ [Clock Cycle Time]. Only instruction count is algorithm-dependent. In the proof-of-concept, we use synthetic noise modeling to simulate measurement variability. Real validation would use Linux \texttt{perf stat -e instructions} to access hardware performance counters.

\textbf{Proxy 3: PR-Style Score (Conformity to Accepted Code Patterns).} This metric would be trained on SWE-bench data using CodeBERT fine-tuned on PR diff acceptance outcomes. In the proof-of-concept, this is implemented as a placeholder (random score generator) because training infrastructure exceeds PoC scope.

\#\#\#\# Data Generation Protocol

\textbf{Dataset:} We selected 50 problems from HumanEval using stratified sampling, supplemented with 50 controlled complexity tasks (O(n), O(n log n), O($n^{2}))$ to test Cohen's d.

\textbf{Solution Generation:} For proof-of-concept, we generated 500 synthetic solutions. Real validation would use CodeLlama-7B-Instruct with temperature=0.8, top\_p=0.95, generating 10 diverse solutions per problem.

\textbf{Repeated Measurements:} Each solution is measured 5 times per proxy metric, yielding 7,500 total evaluations (500 solutions $\times$ 5 repetitions $\times$ 3 proxies).

\textbf{Cross-Platform Simulation:} The proof-of-concept simulates cross-platform measurements by adding Gaussian noise to model hardware-dependent variance.

\#\#\#\# Statistical Analysis

\textbf{Coefficient of Variation (CV):} For each solution and metric, CV = ($\sigma / \mu ) \times$ 100\%, averaged across all solutions. Threshold: CV $\leq 5$\%.

\textbf{Cohen's d (Effect Size):} Using controlled complexity tasks, d = |$\mu _1 - \mu _{2}| / \sigma _pooled$. Threshold: d $\geq 0.8$.

\textbf{Spearman Rank Correlation ($\rho$):} Rank solutions by metric score on each platform, compute Spearman's correlation. Threshold: $\rho \geq 0.8$.

\textbf{Multi-Criteria Gate Validation:} A proxy passes Stage 1 if and only if all three criteria pass. The scoped gate succeeds if $\geq 1$ proxy validates.

\subsubsection{Key Design Decisions}

\textbf{Decision 1: CPU Instruction Count over Wall-Clock Time.} Hardware performance counters provide more stable efficiency measurements by isolating instruction count from system noise. The trade-off: requires Linux perf access.

\textbf{Decision 2: Scoped MUST\_WORK Gate.} We use $\geq 1$ proxy passing as the success criterion rather than requiring all three. This treats partial validation as scientifically valid.

\textbf{Decision 3: Proof-of-Concept with Synthetic Data.} Stage 1 validation uses synthetic measurements before infrastructure investment. This two-phase approach validates methodology before committing resources. The validated methodology transfers; numerical claims require real data confirmation.
